% ABOUTME: Literature review (prior pipelines, open RISC-V cores, branch handling) and background.
% ABOUTME: Owner: Navodit. Background covers CPI, IPC, the iron law and the three kinds of hazard.
\section{Literature Review and Background}
\label{sec:background}

\subsection{Pipelined processors}
The Stanford MIPS processor \cite{hennessy1982} showed that a simple
load/store instruction set lets a single chip run a deep pipeline at a high
clock. MIPS left hazards to the compiler, which had to schedule instructions
so that no result was used too early. Later designs, including ours, detect
hazards in hardware instead. The five-stage pipeline we build (IF, ID, EX,
MEM, WB) is the one Patterson and Hennessy use to teach pipelining
\cite{patterson2021}, together with the forwarding unit, the load-use stall
and the branch flush. Hennessy and Patterson \cite{hennessy2019} give the
iron law of performance that we use to explain every result.

\subsection{RISC-V and open-source cores}
RISC-V is a free and open instruction set from UC Berkeley. Asanovi\'c and
Patterson \cite{asanovic2014} argue that a free ISA lets anyone build and
share implementations without a licence, and that a simple RISC design is the
right base for one. Open RISC-V cores now cover a wide range of designs.
PicoRV32 \cite{picorv32} is optimised for size and clock speed, not
throughput: it takes several cycles per instruction (an average CPI of about
4.2 on Dhrystone). Ibex \cite{ibex} has a two-stage pipeline, fetch and a
combined decode and execute stage. Rocket \cite{rocketchip} is a five-stage
in-order core like ours, with caches and branch prediction, and has been
manufactured and has booted Linux. Our design is the same five stages as
Rocket with none of those extras, so the cost of each hazard can be measured
on its own.

\subsection{Branch handling}
Smith \cite{smith1981} compared branch prediction strategies and showed that
dynamic prediction, which learns from each branch's history, beats fixed
rules. We use the simplest fixed rule, predict not taken, and Section
\ref{sec:results} measures what it costs. That cost is the starting point for
our Phase II work.

\subsection{Performance measures}
For a program of $N$ instructions that takes $C$ clock cycles at clock
frequency $f$,
\begin{equation}
\mathrm{CPI} = \frac{C}{N}, \qquad \mathrm{IPC} = \frac{N}{C}, \qquad
T = \frac{N \cdot \mathrm{CPI}}{f}.
\label{eq:ironlaw}
\end{equation}
The last expression is the iron law of processor performance
\cite{hennessy2019}. Throughput is $f \cdot \mathrm{IPC}$ instructions per
second. A single-cycle CPU has $\mathrm{CPI}=1$ exactly; a pipeline aims for
CPI close to 1 at a much higher $f$.

For the same program, the speedup of a pipeline over the single-cycle design
follows from (\ref{eq:ironlaw}):
\begin{equation}
S = \frac{T_{\text{single}}}{T_{\text{pipe}}} =
\frac{f_{\text{pipe}} / f_{\text{single}}}{\mathrm{CPI}_{\text{pipe}}}.
\label{eq:speedup}
\end{equation}
The pipeline is faster only when its CPI is below its clock ratio.

\subsection{Pipeline cycle accounting}
A five-stage pipeline needs four cycles to fill before the first instruction
completes. After that, every lost cycle has one of two causes in our design:
a stall, where the instruction in ID waits for an operand, or a flush, where
a taken branch discards the two younger instructions already fetched. So
\begin{equation}
C = N + 4 + \text{stalls} + 2 \cdot \text{taken branches}.
\label{eq:cycles}
\end{equation}
Our test suite checks this identity on every directed program and benchmark in both pipeline modes.

\subsection{Hazards}
\emph{Structural hazards} arise when two stages need one resource. Separate
instruction and data memories (a Harvard layout) and a register file with two
read ports and one write port remove them.
\emph{Data hazards} arise when an instruction reads a register that an older
instruction has not yet written. A pipeline either stalls until the value is
in the register file, or forwards it from a later pipeline register straight
to the ALU inputs. A load is the exception: its data only exists after MEM,
so a use in the very next instruction must wait one cycle even with
forwarding (the load-use hazard).
\emph{Control hazards} arise because the next PC after a branch is known only
when the branch resolves. We predict not taken and resolve branches in EX, so
each taken branch costs two cycles.
